\chapter[Höchster Beamter der Landeskirche]{Der höchste Beamte der preußischen Landeskirche}
\quellseite{8}
der Präsident des Berliner Oberkirchenrathes, Dr. Barthhausen, ein Jurist, ist vor einigen Monaten gestorben. Zu seinem Nachfolger ist nun wieder ein Jurist, nicht ein Theologe, ernannt worden, der bisherige Präsident des Consistoriums in Hannover, Vogtke, der der lutherischen Kirche seiner Heimath angehört, nun aber ohne Weiteres in die unirte Landeskirche Preußens übertritt. Es war zwar der Vorschlag gemacht worden, daß ein Theologe an die Spitze berufen werden möchte, in der Hoffnung, daß dadurch der Nothstand der Kirche gebessert werden könnte. Aber dieser Vorschlag ging nicht durch und wurde selbst von kirchlichen Zeitschriften abgelehnt. Eine derselben, die „Kirchliche Wochenschrift“, begründet dies damit, daß eben die modernen Theologen von der Schriftwahrheit abgefallen seien. Das sei noch viel schlimmer, als wenn die Kirche von Juristen geleitet werde. Das Blatt führt aus, daß der Artikel von der Rechtfertigung aus dem Glauben von den modernen Theologen nicht mehr angenommen werde, und was die Stellung zur heiligen Schrift anlange, so hätten viele Theologen bloß noch die Einbanddeckel der Bibel in Händen, aber den Inhalt glaubten sie nicht. Und das sei es, was die Kirche vor allem bei dem Volke in Mißachtung bringe. — Das ist eine scharfe, vernichtende, aber nur zu sehr berechtigte Verurtheilung der modernen Theologie. Und es ist fürwahr kein Wunder, daß große Massen des Volks der Kirche entfremdet werden, wenn diejenigen, die die Lehrer des Volks sein sollten, vom rechten Kirchenglauben abgefallen sind. Davor behüte uns der liebe himmlische Vater und erhalte unsere Kirche bei seinem Wort und Glauben! L. F.
